\scanpage{283}
\mspage{13}

on en déduirait $\pi_0(X)\longrightarrow\varnothing$, \unsure{ou|or|i.e.} $\pi_0(X)=\varnothing$, absurde !.

Reste le cas où $X$ est non vide. Cela résulte du \deleted{fait \unsure{général} \unsure{telle} $X$ \unsure{ayant} une section}\inserted{fait général}

\uline{Lemme} \unsure{2|?}.\unsure{13|14}. Soit \deleted{$\unsure{C}$}$M$ une catégorie ayant un objet initial \uline{strict} $\varnothing_M$ (\deleted{si \unsure{le}} $\operatorname{Hom}_M(X,\varnothing_M)\ne\varnothing$, on a $X=\varnothing_M$), \inserted{un objet final $e_M$,} et telle que pour tt $X\ne\varnothing_M$, $X\longrightarrow e$ ait une section i.e. $\operatorname{Hom}(e,X)\ne\varnothing$.

\deleted{Alors} Soit $W\subset\mathrm{Fl}\,C$, \deleted{Alors tel que \unsure{pour}} \inserted{\textcircled{1\textsuperscript{o}}} ne contenant pas de flèches $\varnothing_M\longrightarrow X$, avec $X\ne\varnothing_M$, et \inserted{\textcircled{2\textsuperscript{o}}} tel que

% margin note written sideways
\inserted{pour tout diagramme
\diagram{\begin{tikzcd}[ampersand replacement=\&]
X_1 \& X_2 \arrow[l,"u_2"']\\
e \arrow[u,"u_1"] \arrow[ur,dashed,"s"'] \&
\end{tikzcd}}
\unsure{avec|où} $u_2\in W$ \deleted{\diagram{\begin{tikzcd}[ampersand replacement=\&]
\unsure{X_1} \& \unsure{X_2} \arrow[l,"{\unsure{u_2}}"']\\
\unsure{e} \arrow[u,"{\unsure{u_1}}"] \arrow[ur,dotted,"{\unsure{s}}"'] \&
\end{tikzcd}} \unsure{avec}} $\exists\,s$ \uline{dans $M$} \inserted{tel} que $u_2s\unsure{=|\simeq}u_1$ dans $MW^{-1}$.}

a) $\varnothing$ est un objet initial \uline{strict} de $MW^{-1}$

b) Pour tt $X\ne\varnothing$, $\operatorname{Hom}_M(\deleted{\unsure{X}}e,\deleted{\unsure{e}}X)\unsure{{-}|\longrightarrow}\operatorname{Hom}_W(e,X)$ est surjectif.

\uline{Dém} Que $\varnothing_M$ soit objet initial de $MW^{-1}$ a été vu (dernière assertion du lemme 2.5, forme duale), prouvons qu'il est strict, i.e. que si \deleted{$\unsure{X\to\varnothing_M}$ \unsure{dans}} $\operatorname{Hom}_W(X,\varnothing)\ne\varnothing$, \inserted{on a $X=\varnothing_M$.} En effet, \deleted{\unsure{sinon}} si on avait $u:X{-}\varnothing$ dans $MW^{-1}$, $u$ \deleted{\unsure{ser}}serait représenté par

\diagram{\begin{tikzcd}[ampersand replacement=\&,column sep=small]
X=X_0 \& X_1 \arrow[l,"W"'] \arrow[r] \& X_2 \& \cdots \arrow[l,"W"'] \& X_{2n-2}\arrow[r] \& X_{2\unsure{n}-1} \& X_{2n}=\varnothing_M \arrow[l,"W"']
\end{tikzcd}}

On procède par récurrence sur la \inserted{demi-}longueur de cette chaîne pour voir que $X=\varnothing$. En effet, l'hyp. sur $W$ implique que $X_{2n-1}=\varnothing_M$, puis celle que $\varnothing_M$ est strict implique $X_{\unsure{2-n|2n-2}}=\varnothing_M$, d'où la conclusion.

% archive sheet number: 282
